% =============================================================================
\bloco{2}{Harder book items and conceptual classification}{Kurlat \S\S1.1, 2.2, 4.3, 5.2, 5.4; Jones ch.~2 \,\textbullet\, 35 points}

\questao{(10 points) \textit{Source: Kurlat, P. (2020), \emph{A Course in Modern Macroeconomics}, Exercise 5.8 ``National Accounts and
The Golden Rule'', pp.~97--98 --- statement reproduced verbatim, all items (a)--(d).}}
\par\smallskip\noindent\fbox{\parbox{\dimexpr\linewidth-2\fboxsep-2\fboxrule}{\small\textbf{Context from the book (added by this list, from
Kurlat \S1.1, p.~21; \S\S4.1--4.4, pp.~55--67; \S\S5.1--5.2, pp.~76--77):} \emph{Accounts:} the conventions in the box of Question~1, plus
depreciation: machines and buildings lose value over time; GDP is \emph{gross} because it is measured before depreciation, so when a firm
subtracts depreciation to compute its profits, the income approach adds it back as a separate line (Ex.~1.10, p.~21). The savings rate is
$S/Y=(Y-C)/Y$, savings as a fraction of income, and in a closed economy savings equal investment, $S=I$ (p.~56).
\emph{Solow model:} with $Y$ output, $K$ capital, $L$ labor (everyone works, so $L$ is also the population), the Cobb--Douglas
production function is $Y=K^{\alpha}L^{1-\alpha}$, $\alpha\in(0,1)$ \,(4.1.2), p.~55 (writing it $AK^{\alpha}L^{1-\alpha}$, with $A$ the
level of technology, ``doesn't make much difference'', p.~77). The population grows at rate $n$, $L_{t+1}=(1+n)L_t$ (p.~55); people
save a constant fraction $s$ of income, so $I=sY$ \,(4.1.3); every piece of capital loses a fraction $\delta$ of its value each period,
so $K_{t+1}=(1-\delta)K_t+I_t$ \,(4.1.4) (pp.~56--57). Per worker, $y\equiv Y/L$, $k\equiv K/L$ and $y=f(k)\equiv F(K/L,1)$ \,(4.2.1)
(p.~57). In a steady state ($ss$), $sf(k_{ss})=(\delta+n)k_{ss}$ \,(4.3.2) and consumption is $c_{ss}=(1-s)y_{ss}$ \,(4.3.1); the
\emph{Golden Rule} is the saving behavior that maximizes steady-state consumption, and the Golden Rule capital stock $k_{gr}$ satisfies
$f'(k_{gr})=\delta+n$ \,(4.3.4) (pp.~62--63). \emph{Factor markets} are perfectly competitive, so the rental rate of capital and the wage
equal the marginal products, $r^K=F_K(K,L)$ \,(4.4.8) and $w=F_L(K,L)$ \,(4.4.9) (p.~66), and all output is paid to workers or to
capital owners, $r^KK+wL=F(K,L)$ (Prop.~4.2, p.~67). With Cobb--Douglas the labor share is $wL/Y=1-\alpha$ \,(5.1.1), so the capital
share is $\alpha$ (p.~76).}}
\par\smallskip
\textit{What this trains:} the three approaches with an investment good and depreciation, then reading factor shares and
the saving rate off the accounts and using them in the golden rule --- Prova~01 Q1 and Q2(b) in one exercise.

\noindent In 2018, the following events happened in Aurum, which is a closed economy that uses the Denarius (plural: Denarii) as its currency.
\begin{itemize}[leftmargin=2.2em,topsep=2pt,itemsep=1pt]
\item Grano, Inc.\ hired Cornelia to plant and harvest 500 tons of wheat on a large plot of land using a tractor. Both the land and
the tractor are owned by Grano, Inc.
  \begin{itemize}[label=--,topsep=0pt,itemsep=0pt]
  \item It paid Cornelia 200 Denarii for her work on the harvest.
  \item It sold the wheat to Panem, Inc.\ for 1,000 Denarii.
  \end{itemize}
\item The tractor the Cornelia used:
  \begin{itemize}[label=--,topsep=0pt,itemsep=0pt]
  \item Was new at the beginning of the year. It was worth 2,000 Denarii.
  \item Was no longer new at the end of the year. By then it was worth 1,800 Denarii.
  \end{itemize}
\item Panem, Inc.\ hired Gaius to grind the wheat in its mill and produce a million loaves of bread.
  \begin{itemize}[label=--,topsep=0pt,itemsep=0pt]
  \item It paid Gaius 1,400 Denarii for his work.
  \item It sold the bread for 4,000 Denarii to the hungry citizens, who eat it.
  \end{itemize}
\item Fabrica, Inc.\ hired Livia to build a combine harvester.
  \begin{itemize}[label=--,topsep=0pt,itemsep=0pt]
  \item It paid Livia 400 Denarii for her work.
  \item It sold the combine harvester to Grano, Inc.\ for 1,000 Denarii.
  \end{itemize}
\end{itemize}
\begin{enumerate}[label=(\alph*),leftmargin=2.2em,topsep=2pt,itemsep=1pt]
\item Construct GDP accounts (in Denarii) for Aurum by production, income and expenditure.
\item What were the income shares of labor and capital?
\item What was the savings rate?
\item Assume the following:
  \begin{itemize}[topsep=0pt,itemsep=0pt]
  \item the savings rate is constant,
  \item the production function has the form:
  \[
  Y=AK^\alpha L^{1-\alpha}
  \]
  and there is no technological progress,
  \item the depreciation rate for the tractor is typical for the capital stock as a whole,
  \item the population is constant.
  \end{itemize}
In the long run, would the level of consumption be higher if the economy slightly increased its saving rate?
\end{enumerate}
\resolespaco
\begin{sol}
\begin{notacao}
\item Grano (G), Panem (P), Fabrica (F): the three firms. $VA_j$, $\Pi_j$ and the identities \eqref{q1:G1}--\eqref{q1:G4}: as in Question~1.
\item $D\ge0$: depreciation of Grano's tractor. $S$: national saving. $Y>0$: GDP; $K>0$: the capital stock; $L>0$: labour.
\item $A>0$: total factor productivity. $\alpha\in(0,1)$: the capital exponent of $Y=AK^{\alpha}L^{1-\alpha}$, equal to the capital
share under competitive factor pay. $\delta\in(0,1)$: depreciation rate. $n\ge0$: labour-force growth (here $n=0$).
\item $s\in(0,1)$: saving rate. $k\equiv K/L$, $y\equiv Y/L$, $c$: capital, output and consumption per worker; the subscript $ss$
marks the steady state. $s_{GR}$: the golden-rule saving rate. $MPK$: marginal product of capital.
\end{notacao}

\passo{(a) accounts}{Grano uses no intermediate inputs (the tractor and the land are capital); the wheat is intermediate for Panem;
the combine harvester is a final good, investment by Grano. By \eqref{q1:G2}:}
\begin{align}
VA_G &= (1{,}000)-(0) \nonumber\\
&= 1{,}000, \label{q3:vaG}\\
VA_P &= (4{,}000)-(1{,}000) \nonumber\\
&= 3{,}000, \label{q3:vaP}\\
VA_F &= (1{,}000)-(0) \nonumber\\
&= 1{,}000. \label{q3:vaF}
\end{align}
Summing \eqref{q3:vaG}--\eqref{q3:vaF} with \eqref{q1:G1}:
\begin{align}
Y &= (1{,}000)+(3{,}000)+(1{,}000) \nonumber\\
&= 5{,}000. \label{q3:Y}
\end{align}
The tractor loses value from 2{,}000 to 1{,}800:
\begin{align}
D &= (2{,}000)-(1{,}800) \nonumber\\
&= 200. \label{q3:D}
\end{align}
Capital income by \eqref{q1:G3}, value added minus wages:
\begin{align*}
\Pi_G &= (1{,}000)-(200)\\
&= 800,\\
\Pi_P &= (4{,}000)-(1{,}000)-(1{,}400)\\
&= 1{,}600,\\
\Pi_F &= (1{,}000)-(400)\\
&= 600.
\end{align*}
Grano's 800 splits into the depreciation \eqref{q3:D} and a profit of
\begin{align*}
\Pi_G-D &= (800)-(200)\\
&= 600 .
\end{align*}
Adding up wages and profits:
\begin{align}
\text{wages} &= (200)+(1{,}400)+(400) \nonumber\\
&= 2{,}000, \label{q3:W}\\
\text{profits} &= (600)+(1{,}600)+(600) \nonumber\\
&= 2{,}800. \label{q3:Pi}
\end{align}
Income: by \eqref{q1:G3}, adding \eqref{q3:W}, \eqref{q3:Pi} and \eqref{q3:D},
\begin{align*}
Y &= (2{,}000)+(2{,}800)+(200)\\
&= 5{,}000 .
\end{align*}
Expenditure: by \eqref{q1:G4}, with $C=4{,}000$ (Panem's food) and $I=1{,}000$ (the harvester),
\begin{align*}
Y &= (4{,}000)+(1{,}000)\\
&= 5{,}000 .
\end{align*}
All three equal \eqref{q3:Y}.
\begin{center}\footnotesize\begin{tabular}{@{}lr@{\qquad}lr@{\qquad}lr@{}}\toprule
\multicolumn{2}{@{}c}{Production} & \multicolumn{2}{c}{Income} & \multicolumn{2}{c@{}}{Expenditure}\\\midrule
Agriculture (Grano) & 1{,}000 & Wages $200+1{,}400+400$ & 2{,}000 & Consumption & 4{,}000\\
Food (Panem) & 3{,}000 & Profits $600+1{,}600+600$ & 2{,}800 & Investment & 1{,}000\\
Machinery (Fabrica) & 1{,}000 & Depreciation & 200 & & \\\midrule
Total & 5{,}000 & Total & 5{,}000 & Total & 5{,}000\\\bottomrule\end{tabular}\end{center}
\passo{(b) shares}{Dividing \eqref{q3:W} by \eqref{q3:Y}:}
\begin{align*}
\text{labor share} &= \frac{(2{,}000)}{(5{,}000)}\\
&= 0.4 .
\end{align*}
Dividing the sum of \eqref{q3:Pi} and \eqref{q3:D} by \eqref{q3:Y} (gross of depreciation, as GDP is gross):
\begin{align}
\text{capital share} &= \frac{(2{,}800)+(200)}{(5{,}000)} \nonumber\\
&= 0.6. \label{q3:ksh}
\end{align}
\passo{(c) saving rate}{Closed economy, no government, so saving is output not consumed, and it equals investment:}
\begin{align*}
S &= Y-C\\
&= (5{,}000)-(4{,}000)\\
&= 1{,}000=I .
\end{align*}
Dividing by \eqref{q3:Y}:
\begin{align}
s &= \frac{(1{,}000)}{(5{,}000)} \nonumber\\
&= 0.2. \label{q3:s}
\end{align}
\passo{(d) golden rule}{With $Y=AK^\alpha L^{1-\alpha}$, competitive factor pay makes the capital share equal to $\alpha$, so by
\eqref{q3:ksh} $\alpha=0.6$. Dividing \eqref{q3:D} by the tractor's value:}
\begin{align}
\delta &= \frac{(200)}{(2{,}000)} \nonumber\\
&= 0.10, \label{q3:delta}
\end{align}
and $n=0$. In steady state saving replaces depreciation, $sY=\delta K$; per worker, substituting $y=Ak^{\alpha}$:
\begin{align}
sy &= \delta k, \nonumber\\
s\big(Ak^{\alpha}\big) &= \delta k. \label{q3:ss}
\end{align}
Divide \eqref{q3:ss} by $\delta k^{\alpha}$ and raise to the power $1/(1-\alpha)$:
\begin{align}
k^{1-\alpha} &= \frac{sA}{\delta}, \nonumber\\
k_{ss} &= \Big(\frac{sA}{\delta}\Big)^{\frac{1}{1-\alpha}}. \label{q3:kss}
\end{align}
Substitute \eqref{q3:kss} into $y=Ak^{\alpha}$:
\begin{align}
y_{ss} &= A\bigg(\Big(\frac{sA}{\delta}\Big)^{\frac{1}{1-\alpha}}\bigg)^{\alpha} \nonumber\\
&= A\Big(\frac{sA}{\delta}\Big)^{\frac{\alpha}{1-\alpha}} \nonumber\\
&= A^{\frac{1}{1-\alpha}}\Big(\frac{s}{\delta}\Big)^{\frac{\alpha}{1-\alpha}}. \label{q3:yss}
\end{align}
Consumption is the part of output not saved:
\begin{equation}
c_{ss}=(1-s)\,y_{ss}. \label{q3:css}
\end{equation}
Substitute \eqref{q3:yss} into \eqref{q3:css}, take logs and differentiate with respect to $s$:
\begin{align}
c_{ss} &= (1-s)\Big(A^{\frac{1}{1-\alpha}}\Big(\frac{s}{\delta}\Big)^{\frac{\alpha}{1-\alpha}}\Big), \nonumber\\
\ln c_{ss} &= \ln(1-s)+\frac{1}{1-\alpha}\ln A+\frac{\alpha}{1-\alpha}\big(\ln s-\ln\delta\big), \nonumber\\
\frac{d\ln c_{ss}}{ds} &= -\frac{1}{1-s}+\frac{\alpha}{(1-\alpha)s}. \label{q3:dlnc}
\end{align}
Evaluate \eqref{q3:dlnc} at $\alpha=0.6$ and \eqref{q3:s}:
\begin{align*}
\frac{d\ln c_{ss}}{ds} &= -\frac{1}{1-(0.2)}+\frac{(0.6)}{\big(1-(0.6)\big)(0.2)}\\
&= -\frac{1}{0.8}+\frac{0.6}{0.4\times0.2}\\
&= -1.25+7.5\\
&= 6.25>0 .
\end{align*}
Multiplying \eqref{q3:dlnc} by $(1-\alpha)s(1-s)>0$ leaves $\alpha(1-s)-(1-\alpha)s=\alpha-s$, so the sign is positive exactly
when $s<\alpha$:
\begin{equation}
s_{GR}=\alpha=0.6. \label{q3:sgr}
\end{equation}
Cross-check with Kurlat's condition $f'(k_{gr})=\delta+n$. From \eqref{q3:ss}, with \eqref{q3:s} and \eqref{q3:delta}:
\begin{align*}
\frac{K}{Y} &= \frac{s}{\delta}\\
&= \frac{(0.2)}{(0.1)}\\
&= 2 .
\end{align*}
So, by the marginal product of the Cobb--Douglas function,
\begin{align*}
MPK &= \alpha\frac{Y}{K}\\
&= \frac{(0.6)}{(2)}\\
&= 0.3>\delta+n=0.1 ;
\end{align*}
capital is below the golden rule. \textbf{Yes:} a slightly higher saving rate raises long-run consumption, by \eqref{q3:sgr}.
\intuicao{The national accounts hand you the two numbers the golden rule needs: the capital share ($\alpha$) and the saving rate ($I/Y$).
Saving 20\% when capital earns 60\% of income means the marginal unit of capital pays far more than it costs to maintain.}
\refbib{Kurlat (2020), \S1.1, Ex.~1.10 (depreciation added back), p.~21; \S4.3, eq.~(4.3.4), p.~63; \S5.2, p.~78. No book answer exists;
every figure is verified by the check script.}
\rubrica{
\item (a) 4 points: each column 1.3; the combine harvester counted as intermediate (GDP 4{,}000): max.\ 1.5 for (a).
\item (b) 2 points; (c) 1 point.
\item (d) 3 points: $s_{GR}=\alpha$ derived or justified 1.5, numbers and conclusion 1.5. Comparing $s$ with the labor share: max.\ 1.
}
\end{sol}

% -----------------------------------------------------------------------------
\needspace{22\baselineskip}   % keep the source line with its context box
\questao{(12 points) \textit{Source: Kurlat, P. (2020), \emph{A Course in Modern Macroeconomics}, Exercise 5.7 ``GDP Accounting,
Interest Rates and Growth Accounting'', pp.~96--97 --- statement reproduced verbatim, all items (a)--(e).}}
\par\smallskip\noindent\fbox{\parbox{\dimexpr\linewidth-2\fboxsep-2\fboxrule}{\small\textbf{Context from the book (added by this list, from
Kurlat \S4.4, pp.~66 and 69; \S5.4, pp.~86--87):} \emph{Accounts:} imports, government output at cost and transfers as in the box of
Question~1; depreciation, the savings rate, the Cobb--Douglas function and its factor shares (labor $1-\alpha$, capital $\alpha$) as in
the box of Question~3. $F_K$ is the marginal product of capital, $\partial Y/\partial K$. \emph{Interest rate:} in competitive markets
the rental rate of capital is $r^K=F_K(K,L)$ \,(4.4.8) (p.~66); someone who borrows at the real interest rate $r$ to invest in capital
earns $r^K-\delta$, so the loan market clears only if
\[ r_{t+1}=r^K_{t+1}-\delta \tag{4.4.12} \]
(p.~69), with $\delta$ the rate at which capital loses value. \emph{Growth accounting:} with $Y_t=F(K_t,L_t,A_t)$ \,(5.4.1), $A$
technology, and GDP growth $g_Y\equiv Y_{t+1}/Y_t-1$, a first-order Taylor approximation gives (5.4.2): $g_Y\approx$ capital share
$\times$ growth rate of capital $+$ labor share $\times$ growth rate of labor $+$ Solow residual, the shares being $F_KK/F$ and $F_LL/F$
(pp.~86--87). The residual is measured as ``Solow Residual $=$ GDP growth $-$ Capital share $\times$ Capital Growth $-$ Labor Share
$\times$ Labor Growth'' and is also called Total Factor Productivity, TFP (p.~87).}}
\par\smallskip
\textit{What this trains:} an open economy with imported capital, imported consumption, exports, government output at cost and
income taxes; then the same accounts feed the interest rate and a Solow residual.

\noindent Proletaria is a mythical country in Central Asia, where the currency is the ruble. The exchange rate is 10 rubles per euro.

Suppose the production function in Proletaria is given by:
\[
Y=K^\alpha L^{1-\alpha}
\]
\begin{enumerate}[label=(\alph*),leftmargin=2.2em,topsep=2pt,itemsep=1pt]
\item Find expressions for the output-to-capital ratio $\frac YK$ and the marginal product of capital $F_K$ as functions of $K$
and $L$.
\end{enumerate}
During 2014, the following events took place in Proletaria.
\begin{itemize}[leftmargin=2.2em,topsep=2pt,itemsep=1pt]
\item Kapitas, the main manufacturer in the country, imported an industrial welder made in Germany, for which it paid 1,000 euro.
\item By the end of the year, the industrial welder was no longer new. Its estimated value in the resale market was 960 euro.
\item 100 workers worked for Kapitas the entire year making screws and nails, using the new welder. Each of them received wages
for 350 rubles.
\item Each of the Kapitas workers paid 100 rubles in income taxes.
\item The total output of Kapitas consisted of 1 million screws and 1 million nails. All of it was exported to Austria for a total
of 4,000 euro.
\item The government of Proletaria employed one of the 100 workers as Chief of the Secret Police (in addition to his factory job)
to maintain law and order, and paid her a salary of 10,000 rubles.
\item The workers ate beef imported from France, which cost a total of 2,500 euro.
\end{itemize}
\begin{enumerate}[label=(\alph*),leftmargin=2.2em,topsep=2pt,itemsep=1pt,start=2]
\item Construct GDP accounts (in rubles) for Proletaria by production, income and expenditure.
\item What were the income shares of labor and capital?
\item Suppose that we know that the capital stock in Proletaria is 100,000 rubles and that the depreciation rate of the
industrial welder is typical for this country's capital stock. What interest rate should we expect to observe?
\item During 2015, additional investment in Proletaria has exceeded depreciation so that the capital stock now stands at
120,000 rubles. The labor force also grew thanks to immigration, and now consists of 105 workers instead of 100. GDP during 2015
was 55,000 rubles (prices were constant). How much of the growth in GDP between 2014 and 2015 can be attributed to growth in
Total Factor Productivity?
\end{enumerate}
\resolespaco
\begin{sol}
\begin{notacao}
\item $Y>0$: GDP; $K>0$: the capital stock; $L>0$: labour (workers); all in the exercise's units, values in rubles.
\item $\alpha\in(0,1)$: the capital exponent of $Y=K^{\alpha}L^{1-\alpha}$, equal to the capital share under competition.
$F_K$: marginal product of capital. $\delta\in(0,1)$: depreciation rate. $r$: the real interest rate.
\item $g_Y$, $g_K$, $g_L$: growth rates of $Y$, $K$, $L$; the Solow residual is the growth of TFP.
\item $C$, $I$, $G$, $X$, $M$, $\Pi$ and the identities \eqref{q1:G1}--\eqref{q1:G4}: as in Question~1. The exchange rate is
10 rubles per euro.
\end{notacao}

\passo{(a) ratios}{Write the technology as}
\begin{equation}
Y=K^{\alpha}L^{1-\alpha}. \label{q4:Y}
\end{equation}
Divide \eqref{q4:Y} by $K$:
\begin{align*}
\frac YK &= K^{\alpha-1}L^{1-\alpha}\\
&= \Big(\frac KL\Big)^{\alpha-1}.
\end{align*}
Differentiate \eqref{q4:Y} with respect to $K$ and use \eqref{q4:Y} again:
\begin{align}
F_K &= \alpha K^{\alpha-1}L^{1-\alpha} \nonumber\\
&= \alpha\frac{\big(K^{\alpha}L^{1-\alpha}\big)}{K} \nonumber\\
&= \alpha\frac YK. \label{q4:FK}
\end{align}
\passo{(b) accounts, in rubles (10 rubles per euro)}{Convert each euro amount, and compute the wage bill and the taxes:}
\begin{align*}
\text{welder (imported investment)} &= (1{,}000)\times(10)\\
&= 10{,}000,\\
\text{exports} &= (4{,}000)\times(10)\\
&= 40{,}000,\\
\text{beef (imported consumption)} &= (2{,}500)\times(10)\\
&= 25{,}000,\\
\text{Kapitas wages} &= (100)\times(350)\\
&= 35{,}000,\\
\text{income taxes} &= (100)\times(100)\\
&= 10{,}000.
\end{align*}
Kapitas's gross profit, by \eqref{q1:G3}:
\begin{align*}
\Pi &= (40{,}000)-(35{,}000)\\
&= 5{,}000 .
\end{align*}
It splits into depreciation and net profit:
\begin{align*}
\text{depreciation} &= \big((1{,}000)-(960)\big)\times(10)\\
&= 400,\\
\text{net profit} &= (5{,}000)-(400)\\
&= 4{,}600.
\end{align*}
The police chief's salary, 10{,}000, is government output valued at cost. The income taxes are a transfer and appear nowhere.
Product, by \eqref{q1:G1}--\eqref{q1:G2}:
\begin{align}
Y &= (40{,}000)+(10{,}000) \nonumber\\
&= 50{,}000. \label{q4:Y0}
\end{align}
Income, by \eqref{q1:G3}:
\begin{align*}
Y &= (35{,}000)+(10{,}000)+(5{,}000)\\
&= 50{,}000 .
\end{align*}
Imports are the welder plus the beef; expenditure, by \eqref{q1:G4}:
\begin{align*}
M &= (10{,}000)+(25{,}000)\\
&= 35{,}000,\\
Y &= (25{,}000)+(10{,}000)+(10{,}000)+(40{,}000)-(35{,}000)\\
&= 50{,}000 .
\end{align*}
\begin{center}\footnotesize\begin{tabular}{@{}lr@{\qquad}lr@{\qquad}lr@{}}\toprule
\multicolumn{2}{@{}c}{Production} & \multicolumn{2}{c}{Income} & \multicolumn{2}{c@{}}{Expenditure}\\\midrule
Manufacturing & 40{,}000 & Wages (Kapitas) & 35{,}000 & Consumption & 25{,}000\\
Public administration & 10{,}000 & Wages (government) & 10{,}000 & Investment & 10{,}000\\
 & & Profit 4{,}600 $+$ deprec.\ 400 & 5{,}000 & Government & 10{,}000\\
 & & & & Exports & 40{,}000\\
 & & & & Imports (welder $+$ beef) & $-35{,}000$\\\midrule
Total & 50{,}000 & Total & 50{,}000 & Total & 50{,}000\\\bottomrule\end{tabular}\end{center}
\passo{(c) shares}{Dividing all wages by \eqref{q4:Y0}:}
\begin{align*}
\text{labor share} &= \frac{(35{,}000)+(10{,}000)}{(50{,}000)}\\
&= 0.9 .
\end{align*}
Dividing the gross profit by \eqref{q4:Y0}:
\begin{align}
\alpha=\text{capital share} &= \frac{(5{,}000)}{(50{,}000)} \nonumber\\
&= 0.1. \label{q4:alpha}
\end{align}
\passo{(d) interest rate}{The depreciation rate is the value lost (the same 40 euros behind the 400 rubles of (b)) over the
value at the start:}
\begin{align}
\delta &= \frac{(1{,}000)-(960)}{(1{,}000)} \nonumber\\
&= 0.04. \label{q4:delta}
\end{align}
The interest rate is the marginal product of capital net of depreciation; substitute \eqref{q4:FK}:
\begin{align}
r &= F_K-\delta \nonumber\\
&= \Big(\alpha\frac YK\Big)-\delta. \label{q4:r}
\end{align}
Substitute \eqref{q4:alpha}, \eqref{q4:Y0}, $K=100{,}000$ and \eqref{q4:delta} into \eqref{q4:r}:
\begin{align*}
r &= (0.1)\times\frac{(50{,}000)}{(100{,}000)}-(0.04)\\
&= 0.05-0.04\\
&= 0.01,
\end{align*}
i.e.\ 1\%.
\passo{(e) growth accounting}{Growth rates from one year to the next:}
\begin{align*}
g_Y &= \frac{(55{,}000)}{(50{,}000)}-1\\
&= 10\%,\\
g_K &= \frac{(120{,}000)}{(100{,}000)}-1\\
&= 20\%,\\
g_L &= \frac{(105)}{(100)}-1\\
&= 5\%.
\end{align*}
The Solow residual is output growth not explained by factor growth, weighted by the shares:
\begin{equation}
\text{residual}=g_Y-\alpha\,g_K-(1-\alpha)\,g_L. \label{q4:res}
\end{equation}
Substitute the growth rates and \eqref{q4:alpha} into \eqref{q4:res}:
\begin{align*}
\text{residual} &= (10\%)-(0.1)\times(20\%)-\big(1-(0.1)\big)\times(5\%)\\
&= 10\%-0.1\times20\%-0.9\times5\%\\
&= 10\%-2\%-4.5\%\\
&= 3.5\%.
\end{align*}
So 3.5 of the 10 percentage points, 35\% of the growth, is TFP. With exact logs in \eqref{q4:res}:
\begin{align*}
\text{residual} &= (\ln1.1)-(0.1)(\ln1.2)-(0.9)(\ln1.05)\\
&= 0.0953-0.0182-0.0439\\
&= 0.0332,
\end{align*}
34.8\% of $\ln1.1$.
\intuicao{Imports are subtracted because $C$ and $I$ include goods made abroad; the welder raises investment and imports by the same
amount and leaves GDP untouched. The income side is what gives $\alpha$, and $\alpha$ is what turns GDP into $r$ and into TFP.}
\refbib{Kurlat (2020), \S1.1, Exs.~1.6--1.8, pp.~19--20; \S5.2, $r=F_K-\delta$, p.~78; \S5.4, eq.~(5.4.2), p.~87. No book answer exists.}
\rubrica{
\item (a) 1 point; (b) 5 points (each column 1.5; conversion at 10 rubles per euro 0.5).
\item Income taxes counted in GDP, or the welder counted without its import: max.\ 2 for (b).
\item (c) 2 points (government wages are labor income; using $35{,}000/40{,}000$ with a stated reason: 1.5).
\item (d) 2 points: $r$ without $-\delta$: max.\ 1. (e) 2 points.
}
\end{sol}

% -----------------------------------------------------------------------------
\questao{(10 points) Conceptual --- answer in words, no numbers. For each event, say whether it adds to this year's GDP;
if it does, name the producer whose value added rises, the income it generates and the final use ($C$, $I$, $G$ or $X$) it
enters; if it does not, say why. Two to four sentences each.}
\noindent\textbf{a)} A search engine is free for households and is paid for by advertising sold to firms.\\
\textbf{b)} A bank's own IT staff write a risk-management program the bank will use for years; a drug company's
laboratory runs a research project for its own future products.\\
\textbf{c)} A household sells its five-year-old car to another household through a used-car dealer.\\
\textbf{d)} One family owns the house it lives in; next door an identical house is rented to a tenant.\\
\textbf{e)} An accounting firm prepares one tax return for a manufacturing company and one for a household.\\
\textbf{f)} The government pays retirement pensions and pays the salaries of public-school teachers.
\resolespaco
\begin{sol}
\begin{notacao}
\item $C$, $I$, $G$, $X$: the final uses of \eqref{q1:G4} (consumption, investment, government purchases, exports).
\item Value added, income and the three approaches: \eqref{q1:G1}--\eqref{q1:G3} of Question~1. No numbers are needed.
\end{notacao}
\passo{a)}{The ads are a service bought by firms, so they are an \emph{intermediate} input: the network's value added is counted, but no
final use appears for it; it is embedded in the price of the advertisers' goods. The search itself has a zero price and is no market
transaction, so it adds nothing beyond that. The users' gain is real welfare that GDP does not measure.}
\passo{b)}{Both are \emph{gross fixed capital formation} in intellectual property products: durable assets that yield services for several
years. With no market price, each is valued at its production cost. Production: the bank's and the company's output includes the
asset; income: the wages paid plus an operating surplus; expenditure: $I$. It is not consumption (no household buys it) and not an
intermediate input (it is not used up within the year).}
\passo{c)}{The car was produced, and counted, five years ago; changing its owner is not new production. Only the dealer's service is
current output: its margin is the dealer's value added, its profit and wages, and the buyer's consumption of a trade service.}
\passo{d)}{Both houses produce housing services that enter $C$. For the tenant, the rent is paid in the market; for the owner, an
\emph{imputed} rent is counted as both consumption and income of the owner. Without it, GDP would fall every time a tenant bought the
house she lives in, with no change in what is produced.}
\passo{e)}{The same service is \emph{intermediate} when the company buys it (a cost of making its goods, already inside their price) and
\emph{final} when the household buys it ($C$). The accounting firm's value added counts in both cases; only the final-use column
differs. The classification follows the buyer and the use, never the product.}
\passo{f)}{Pensions are transfers: no good or service is produced in exchange, so they enter neither $G$ nor GDP. Teachers produce
education with no market price; it is valued at cost: value added of public education $=$ the wages $=$ income, and the final use is
$G$.}
\intuicao{One question settles every case: was something produced this year, and who used it up? Firms using it $\Rightarrow$
intermediate; households $\Rightarrow C$; firms keeping it for years $\Rightarrow I$; government $\Rightarrow G$; nothing produced $\Rightarrow$ 0.}
\refbib{Kurlat (2020), \S1.1: inventories and investment (Ex.~1.3, p.~17), financial investment (Ex.~1.4, p.~18), durables and imputed rent
(Ex.~1.5, pp.~18--19), government at cost and transfers (Exs.~1.7--1.8, p.~20), non-market activity (Ex.~1.11, p.~22); \S2.2, p.~33.
Jones (2020), ch.~2, p.~22: intellectual property products in $I$; government purchases vs.\ transfers.}
\rubrica{
\item 10/6 points per item, judged on the classification and the reason.
\item Own-account software or R\&D called consumption (the Prova~01 Q1(c) error), or the free service counted as a final use: 0 for that item.
\item Right classification with no reason: half the item.
}
\end{sol}

% -----------------------------------------------------------------------------
\needspace{28\baselineskip}   % keep the source line with its context box
\questao{(3 points) \textit{Source: Kurlat, P. (2020), \emph{A Course in Modern Macroeconomics}, Exercise 1.4 ``Drugs and
Prostitution'', p.~29, and Exercise 2.4 ``Construction and the Armed Forces'', p.~41; Jones, C.\,I. (2020), \emph{Macroeconomics},
5th ed., ch.~2, Exercise~9 ``Earthquakes and GDP'', p.~41 --- statements reproduced verbatim; the closing instruction is added
by this list.}}
\par\smallskip\noindent\fbox{\parbox{\dimexpr\linewidth-2\fboxsep-2\fboxrule}{\small\textbf{Context from the books (added by this list,
from Kurlat \S1.1, pp.~15 and 21; Kurlat \S2.2, pp.~33--34; Jones ch.~2, pp.~19, 22 and 26):} \emph{Kurlat 1.4.} GDP records the value
of everything produced in the country in a year (p.~15); activity that involves no market transaction is usually not included: ``Most
of the time, however, we compute the value of an activity only if it is sold in the market'' (p.~21). The article is not reproduced here;
its address names its subject: Italy's inclusion of illicit activities in its GDP figures.
\emph{Kurlat 2.4.} Jones and Klenow measure the living standard of a country, Utilia, by the number $\lambda$ that makes a person
indifferent between living a year as a random resident of Utilia and as a random resident of the US with everyone's consumption
multiplied by $\lambda$ (p.~33), using
\[ u(c,l,a)=\mathbb{E}\Big[\Big(\bar u+\frac{c^{1-\sigma}}{1-\sigma}-\theta(1-l)^2\Big)a\Big], \tag{2.2.1} \]
\[ u\big(\lambda c^{US},l^{US},a^{US}\big)=u\big(c^{Utilia},l^{Utilia},a^{Utilia}\big), \tag{2.2.2} \]
where $c$ is consumption, $l$ leisure (the fraction of time spent on activities not counted in GDP), $a$ equals 1 if alive and 0 if dead,
and $\bar u$, $\sigma$, $\theta$ are parameters (pp.~33--34). What enters (2.2.1) is consumption, not production: output dedicated to
investment or exported gives no utility during the year (investment benefits future residents), and public and private consumption are
assumed to give the same utility, so what enters is $C+G$ of (1.1.1) (box of Question~1), more precisely $C$ plus the government-consumption part of $G$
(p.~34). \emph{Jones 9.} GDP is ``the market value of the final goods and services produced in an economy over a certain period''
(p.~19); the construction of new homes is part of investment (p.~22); GDP counts only new production (p.~26).}}
\par\smallskip
\textit{What this trains:} separating what GDP measures (market production) from what it does not (legality, welfare, the
stock of wealth).

\noindent\textbf{Kurlat 1.4 Drugs and Prostitution.} Read the following article:
\texttt{http://\allowbreak www.economist.com/\allowbreak news/\allowbreak finance-\allowbreak and-\allowbreak economics/\allowbreak
21603073-\allowbreak italys-\allowbreak inclusion-\allowbreak illicit-\allowbreak activities-\allowbreak its-\allowbreak figures-\allowbreak
excites-\allowbreak much-\allowbreak interest-\allowbreak sex}. What do you think? Should drug
production and prostitution be included in the calculation of GDP?

\noindent\textbf{Kurlat 2.4 Construction and the Armed Forces.} Suppose we compared welfare in two otherwise identical countries.
Country A dedicates 50\% of GDP to maintaining a large army. Country B dedicates 50\% of GDP to building houses. In which of the
two countries would welfare be higher, according to the measure used by Jones and Klenow. What do you think of this?

\noindent\textbf{Jones 9. Earthquakes and GDP:} Suppose the rural part of a country is hit by a major earthquake that destroys
10 percent of the country's housing stock. The government and private sector respond with a major construction effort to help
rebuild houses. Discuss how this episode is likely to affect (a) the economic well-being of the people in the country and (b) the
economy's measured GDP.

\noindent\textit{Added (not in the book):} answer each of the three in words.
\resolespaco
\begin{sol}
\begin{notacao}
\item $C$, $G$, $I$: consumption, government purchases and investment, as in \eqref{q1:G4}.
\item Jones--Klenow measure: Kurlat's welfare index, built on consumption $C+G$ (Kurlat 2020, \S2.2).
\end{notacao}
\passo{Kurlat 1.4}{GDP measures market production, not its moral worth; drugs and prostitution are voluntary transactions that pay
incomes and meet final demand, so the accounts' logic says include them. Including them also keeps GDP comparable across
countries with different laws and over time when laws change (legalisation would otherwise raise measured GDP with no change in
activity). Against: the figures are estimates built on strong assumptions. Theft, by contrast, is excluded: it is a forced transfer,
not production (Williamson 2014, solutions manual, ch.~2, problem~9).}
\passo{Kurlat 2.4}{GDP is the same in both. The Jones--Klenow measure, as Kurlat implements it, puts $C+G$ in utility: military spending
is government consumption and counts, while housebuilding is investment and does not. So country A scores higher this year. That
is hard to accept: it assumes a unit of defence is worth a unit of private consumption, and it ignores that the houses will deliver
housing services --- consumption, through rents and imputed rents --- for decades. The measure is a one-year snapshot; investment
pays off to future residents.}
\passo{Jones 9}{(a) Well-being falls: a tenth of the housing stock is lost, and the rebuilding uses labour and materials that would
otherwise have made something new (the broken-window fallacy). (b) Measured GDP probably rises: reconstruction is new residential
investment, and GDP counts the flow of production, not the stock of wealth destroyed. Confirmed: Jones \emph{Instructor's Manual}
(2016), ch.~2, Exercise~8, pp.~12--13.}
\refbib{Kurlat (2020), \S2.2, pp.~33--34 (Jones--Klenow; $C+G$ in utility, p.~34); Jones (2020), ch.~2.}
\rubrica{
\item 1 point per exercise. Any answer that equates GDP with welfare: 0 for that exercise.
\item Kurlat 2.4 answered ``B is better off'' without noting that the measure puts $G$ in and $I$ out: half.
}
\end{sol}
